\documentclass[11pt]{article}

\usepackage[a4paper,margin=1in]{geometry}
\usepackage[T1]{fontenc}
\usepackage{lmodern}
\usepackage{microtype}
\usepackage{parskip}
\usepackage{graphicx}
\usepackage{amsmath,amssymb}
\usepackage{booktabs}
\usepackage{array}
\usepackage{multirow}
\usepackage{caption}
\usepackage{float}
\usepackage{enumitem}
\usepackage{xcolor}
\usepackage{titlesec}
\usepackage{fancyhdr}
\usepackage{listings}
\usepackage[most]{tcolorbox}

\definecolor{primary}{HTML}{1B3A5C}
\definecolor{accent}{HTML}{2E86AB}
\definecolor{accent2}{HTML}{C9622A}
\definecolor{panelbg}{HTML}{F4F6F8}

\renewcommand{\arraystretch}{1.2}

\titleformat{\section}
  {\normalfont\Large\bfseries\color{primary}}
  {\thesection}{0.6em}{}[{\color{accent}\titlerule[1.1pt]}]
\titleformat{\subsection}
  {\normalfont\large\bfseries\color{primary}}
  {\thesubsection}{0.5em}{}
\titleformat{\subsubsection}
  {\normalfont\bfseries\color{accent}}
  {\thesubsubsection}{0.5em}{}
\titlespacing*{\section}{0pt}{1.6em}{0.8em}
\titlespacing*{\subsection}{0pt}{1.2em}{0.5em}

\pagestyle{fancy}
\fancyhf{}
\renewcommand{\headrulewidth}{0.4pt}
\renewcommand{\footrulewidth}{0pt}
\fancyhead[L]{\small\color{primary}\bfseries CS683 PA-1 - Task 2}
\fancyhead[R]{\small\color{primary}\nouppercase{\rightmark}}
\fancyfoot[C]{\small\color{primary}\thepage}

\lstdefinestyle{cmdstyle}{
    backgroundcolor=\color{panelbg},
    basicstyle=\ttfamily\small,
    frame=single,
    rulecolor=\color{accent},
    framerule=0.6pt,
    breaklines=true,
    columns=fullflexible,
    xleftmargin=4pt, xrightmargin=4pt,
    aboveskip=8pt, belowskip=8pt,
}
\lstset{style=cmdstyle}

\newtcolorbox{insightbox}{
    enhanced, breakable, colback=accent!6, colframe=accent,
    boxrule=0.8pt, arc=3pt, left=8pt, right=8pt, top=6pt, bottom=6pt,
    fonttitle=\bfseries, coltitle=white, colbacktitle=accent,
    title={Key Insight}, before skip=8pt, after skip=10pt,
}
\newtcolorbox{todobox}[1][Fill in]{
    enhanced, breakable, colback=accent2!4, colframe=accent2!70!black,
    boxrule=0.6pt, arc=2pt, left=8pt, right=8pt, top=5pt, bottom=5pt,
    fonttitle=\bfseries\small, coltitle=accent2!70!black,
    title={TODO --- #1}, before skip=6pt, after skip=10pt,
}

\usepackage[colorlinks=true,linkcolor=accent,citecolor=accent,urlcolor=accent]{hyperref}

\graphicspath{{../../plots/}}

\begin{document}

\begin{titlepage}
    \centering
    \vspace*{2.2cm}
    {\color{accent}\rule{0.5\textwidth}{1.2pt}}\\[0.6cm]
    {\Huge\bfseries\color{primary} CS683 PA-1}\\[0.4cm]
    {\LARGE\bfseries\color{primary} Task 2 Report}\\[0.3cm]
    {\Large Matrix Multiplication (SGEMM)}\\[0.6cm]
    {\color{accent}\rule{0.5\textwidth}{1.2pt}}\\[1.4cm]

    {\large Software Prefetching \(\cdot\) SIMD \(\cdot\) llama.cpp Integration}\\[2cm]

    \vfill
    {\large \today}
\end{titlepage}

\tableofcontents
\newpage

\section*{Setup}
\addcontentsline{toc}{section}{Setup}

All measurements were taken on a 13th Gen Intel Core i7-13650HX. Every build and every run was
pinned to CPU~0, a P-core, with \texttt{taskset -c 0}. L1-D size is 48KB and line size is 64B.

\section{Task 2A: Software Prefetching}

\subsection{Implementation}

\texttt{matmul\_prefetch} tiles $M$, $N$, and $K$ with block size $T=64$. Within a tile, the $K$-loop is unrolled by 16 and
\texttt{\_mm\_prefetch(\ldots, \_MM\_HINT\_T0)} issues a lookahead load for
both operands every 16 elements, at a chosen distance ahead of use. The
lookahead is bounded by the true row length $K$ (not the current $K$-tile's
edge), so prefetching correctly continues into the next $K$-tile of the same
row for large distances. \texttt{\_mm\_prefetch} was used over
\texttt{\_\_builtin\_prefetch} for direct control over the locality hint.

\subsection{Speedup vs. Matrix Size}

Hint fixed at \texttt{\_MM\_HINT\_T0}, $T=64$, distance in $\{32,64\}$,
$N\in\{128,256,512,1024,1536,2048,3072\}$. Speedup is vs.\
\texttt{matmul\_naive}.

\begin{figure}[H]
    \centering
    \includegraphics[width=0.75\textwidth]{task2a_size_speedup.png}
    \caption{Speedup of \texttt{matmul\_prefetch} over \texttt{matmul\_naive}
    vs.\ matrix size, for distance $=32$ and $=64$.}
    \label{fig:t2-prefetch-vs-size}
\end{figure}

\begin{table}[H]
    \centering
    \begin{tabular}{rrr}
        \toprule
        \textbf{$N$} & \textbf{Speedup (dist$=32$)} & \textbf{Speedup (dist$=64$)} \\
        \midrule
        128  & 1.33x & 1.53x \\
        256  & 1.38x & 1.23x \\
        512  & 1.34x & 1.29x \\
        1024 & 1.35x & 1.38x \\
        1536 & 1.31x & 1.37x \\
        2048 & 1.36x & 1.35x \\
        3072 & 1.39x & 1.44x \\
        \bottomrule
    \end{tabular}
    \caption{Speedup vs.\ \texttt{matmul\_naive} across matrix sizes. Speedup
    is consistently in the $1.2$--$1.5\times$ range with no strong monotonic
    trend; distance $=64$ has the higher mean (1.37x vs.\ 1.35x) and wins at
    the largest size tested ($N=3072$).}
    \label{tab:t2-prefetch-vs-size}
\end{table}

\begin{figure}[H]
    \centering
    \includegraphics[width=0.95\textwidth]{task2a_size_cache_metrics.png}
    \caption{L1D, L2, LLC misses and \texttt{sw\_prefetch\_access.any} vs.\
    matrix size (log-log), for distance $=32$ and $=64$. All four metrics
    grow monotonically with $N^3$-scale work; L2 misses cross into the
    billions and LLC misses into the tens of millions by $N=3072$, past the
    point where the working set fits in LLC.}
    \label{fig:t2-prefetch-vs-size-metrics}
\end{figure}

\textit{Why flat:} both \texttt{matmul\_naive} and \texttt{matmul\_prefetch} read
$A$ and $B$ with unit stride (the harness's NT layout stores $B$ row-major by
output column, so no transpose is needed), so naive is already reasonably
cache-friendly and not pathologically bad at any size. Prefetching only hides
a roughly constant fraction of the memory stall time behind compute at every
$N$ -- it does not change what fraction of the growing miss counts in
Fig.~\ref{fig:t2-prefetch-vs-size-metrics} are avoidable -- so the speedup
stays in a narrow band rather than growing with $N$.

\subsection{Speedup vs. Prefetch Distance}

Hint fixed at \texttt{\_MM\_HINT\_T0}, $T=64$, distance swept in
$\{8,16,32,64,128,256\}$ floats at $N=1024,2048$. Speedup is vs.\
\texttt{matmul\_naive} on the same workload/seed.

\begin{figure}[H]
    \centering
    \includegraphics[width=0.75\textwidth]{task2a_distance_speedup.png}
    \caption{Speedup vs.\ prefetch distance for \texttt{matmul\_prefetch}.}
    \label{fig:t2-prefetch-distance}
\end{figure}

\begin{table}[H]
    \centering
    \begin{tabular}{rrr}
        \toprule
        \textbf{Distance (floats)} & \textbf{Speedup ($N{=}1024$)} & \textbf{Speedup ($N{=}2048$)} \\
        \midrule
        8   & 1.32x & 1.29x \\
        16  & 1.10x & 1.15x \\
        32  & 1.33x & 1.32x \\
        64  & 1.31x & 1.36x \\
        128 & 1.25x & 1.30x \\
        256 & 1.18x & 1.26x \\
        \bottomrule
    \end{tabular}
    \caption{Speedup vs.\ \texttt{matmul\_naive} by prefetch distance. Best mean
    speedup at distance $=64$ (1.33x).}
    \label{tab:t2-prefetch-distance}
\end{table}

\begin{figure}[H]
    \centering
    \includegraphics[width=0.95\textwidth]{task2a_distance_cache_metrics.png}
    \caption{L1D, L2, LLC misses and \texttt{sw\_prefetch\_access.any} vs.\
    prefetch distance. L2/LLC misses rise gradually with distance (prefetched
    lines evicted before use); \texttt{sw\_prefetch\_access.any} falls at
    large distances.}
    \label{fig:t2-prefetch-distance-metrics}
\end{figure}

\textit{Why $64$ wins:} it gives the load pipeline enough lead time to complete
the DRAM/LLC fetch before the data is touched, without prefetching so far
ahead that the line is evicted before use -- L2/LLC misses in
Fig.~\ref{fig:t2-prefetch-distance-metrics} climb again past distance$=64$,
consistent with early prefetches going stale. Distance$=16$ is an outlier low
($1.12\times$); it coincides with the $K$-unroll granularity of 16, so at that
one distance the prefetch and its consuming load land in the same unrolled
step, leaving almost no lead time -- effectively no prefetching at all.

\subsection{Speedup vs. Prefetch Locality Hint}

Distance fixed at $64$ floats, $T=64$, hint swept over
\texttt{\_MM\_HINT\_\{T0,T1,T2,NTA\}} at $N=1024,2048$. Speedup is vs.\
\texttt{matmul\_naive} on the same workload/seed.

\begin{figure}[H]
    \centering
    \includegraphics[width=0.75\textwidth]{task2a_hint_speedup.png}
    \caption{Speedup for \texttt{\_MM\_HINT\_T0 / T1 / T2 / NTA} locality hints.}
    \label{fig:t2-prefetch-hint}
\end{figure}

\begin{table}[H]
    \centering
    \begin{tabular}{rrr}
        \toprule
        \textbf{Hint} & \textbf{Speedup ($N{=}1024$)} & \textbf{Speedup ($N{=}2048$)} \\
        \midrule
        T0  & 1.37x & 1.38x \\
        T1  & 1.37x & 1.33x \\
        T2  & 1.35x & 1.35x \\
        NTA & 1.34x & 1.36x \\
        \bottomrule
    \end{tabular}
    \caption{Speedup vs.\ \texttt{matmul\_naive} by locality hint. Best mean
    speedup at \texttt{T0} (1.38x).}
    \label{tab:t2-prefetch-hint}
\end{table}

\begin{figure}[H]
    \centering
    \includegraphics[width=0.95\textwidth]{task2a_hint_cache_metrics.png}
    \caption{L1D, L2, LLC misses and \texttt{sw\_prefetch\_access.any} vs.\
    locality hint. LLC misses are lowest for \texttt{T1}, highest for
    \texttt{NTA}; L1D misses and \texttt{sw\_prefetch\_access.any} are
    essentially flat across hints.}
    \label{fig:t2-prefetch-hint-metrics}
\end{figure}

\textit{Why T0 (marginally) wins:} with distance already near-optimal, the
hint only decides which cache levels the fetched line lands in, not whether it
arrives in time. \texttt{T0} fills every level including L1D, so the line is
already at the point of use when reused; \texttt{T1}/\texttt{T2}/\texttt{NTA}
place it in L2, LLC, or a non-temporal buffer instead, adding a small extra
hop up to L1D on the actual load. That hop is cheap, so the spread between
hints ($1.34$--$1.38\times$) is much smaller than the spread across distances.

\subsection{Effect of Hardware Prefetchers}

All four L2/L3 hardware prefetchers were disabled via MSR \texttt{0x1A4}
(\texttt{wrmsr -p 0 0x1a4 0xf}, restored to \texttt{0x0} after) at the best
config (distance$=64$, hint$=$T0, $T=64$).

\begin{table}[H]
    \centering
    \begin{tabular}{rrrrr}
        \toprule
        \textbf{$N$} & \textbf{HW prefetch} & \textbf{naive (ms)} & \textbf{prefetch (ms)} & \textbf{speedup} \\
        \midrule
        1024 & ON  & 476.3  & 348.0  & 1.37x \\
        1024 & OFF & 476.7  & 383.7  & 1.24x \\
        2048 & ON  & 4157.1 & 3069.0 & 1.35x \\
        2048 & OFF & 4930.7 & 3367.6 & 1.46x \\
        \bottomrule
    \end{tabular}
    \caption{Effect of disabling all four HW prefetchers. At $N=1024$, both
    kernels are barely affected. At $N=2048$, \texttt{naive} slows down by
    $19$--$44\%$ (a repeat run without \texttt{perf} overhead measured naive
    OFF at $5964.1$ms, i.e.\ up to $44\%$ slower, showing high run-to-run
    variance once HW prefetch is off) while \texttt{prefetch} slows by only
    $\sim\!10\%$ -- so the naive/prefetch speedup ratio \emph{increases} when
    HW prefetch is disabled.}
    \label{tab:t2-hwpf}
\end{table}

\begin{table}[H]
    \centering
    \begin{tabular}{rrrrr}
        \toprule
        \textbf{Metric} & \textbf{naive ON} & \textbf{naive OFF} & \textbf{prefetch ON} & \textbf{prefetch OFF} \\
        \midrule
        L1D miss & 135{,}566{,}760 & 2{,}407{,}962{,}620 & 1{,}003{,}615{,}281 & 5{,}065{,}277{,}174 \\
        L2 miss  & 3{,}743{,}418{,}813 & 3{,}281{,}496{,}014 & 4{,}188{,}874{,}265 & 3{,}376{,}038{,}747 \\
        LLC miss & 2{,}937{,}937 & 731{,}057{,}910 & 5{,}025{,}087 & 729{,}032{,}068 \\
        sw\_prefetch & 136{,}559 & 160{,}125 & 5{,}209{,}576{,}955 & 5{,}208{,}353{,}385 \\
        \bottomrule
    \end{tabular}
    \caption{Cache metrics at $N=2048$, HW prefetch ON vs.\ OFF. LLC misses
    jump $\sim\!250\times$ for both kernels when HW prefetch is disabled
    (2.9M$\to$731M naive, 5M$\to$729M prefetch) -- with the streaming
    prefetcher gone, far more L2 misses actually reach the LLC/DRAM instead
    of being hidden. \texttt{sw\_prefetch\_access.any} is unchanged (software
    prefetch issue rate does not depend on the HW prefetcher setting).}
    \label{tab:t2-hwpf-cache}
\end{table}

\subsection{Questions}
\begin{enumerate}[leftmargin=1.6em, itemsep=0.4em]
    \item \textbf{Trend in speedup with matrix size:} Roughly flat, $1.2$--$1.5\times$
    from $N=128$ to $3072$, with no strong monotonic trend (Table~\ref{tab:t2-prefetch-vs-size}).
    This is because \texttt{matmul\_naive} is already unit-stride on both
    operands, so it is never
    that slow to begin with. Prefetching only removes a roughly
    constant fraction of stall cycles at every size, rather something
    that scales uniformly with $N$. The raw miss counts in
    Fig.~\ref{fig:t2-prefetch-vs-size-metrics} grow by orders of magnitude
    with $N$, but that growth is common to both kernels, so the \emph{ratio}
    stays flat even though the \emph{absolute} misses avoided keep growing.

    \item \textbf{Best prefetch distance:} $64$ floats, mean speedup $1.33\times$
    across $N{=}1024,2048$ (Table~\ref{tab:t2-prefetch-distance}), narrowly
    ahead of $32$ ($1.33\times$) and clearly ahead of $128$/$256$ ($1.27\times$/$1.22\times$).
    $64$ floats gives the load pipeline enough lead time to complete the
    DRAM/LLC fetch before the data is touched; going further ahead (128, 256)
    starts letting the line get evicted again before use, which is exactly
    why L2/LLC misses climb back up past distance$=64$ in
    Fig.~\ref{fig:t2-prefetch-distance-metrics}. Distance $16$ is an outlier
    low ($1.12\times$). It may be because it coincides with the kernel's $K$-unroll
    granularity of 16: the prefetch and the load it is meant to hide land in
    the same unrolled step, leaving almost no lead time. But it also might be due to CPU core's non deterministic clock cycles.

    \item \textbf{Cache fill level (locality hint) giving maximum speedup:}
    \texttt{\_MM\_HINT\_T0} (fill all cache levels), mean speedup $1.38\times$
    across $N{=}1024,2048$; \texttt{T1}, \texttt{T2}, \texttt{NTA} all $1.35\times$
    (Table~\ref{tab:t2-prefetch-hint}). With distance already tuned, the hint
    only decides which levels the fetched line lands in, not whether it
    arrives in time: \texttt{T0} places it directly in L1D, the level it will
    actually be read from, while \texttt{T1}/\texttt{T2}/\texttt{NTA} land it
    in L2, LLC, or a non-temporal buffer and pay a small extra hop up to L1D
    on the real load. That hop is cheap relative to a full memory fetch, so
    the spread across hints ($1.34$--$1.38\times$) is much narrower than the
    spread across distances ($1.12$--$1.33\times$).

    \item \textbf{L1D / L2 / LLC miss trends across sizes and parameters:}
    All three miss counts grow monotonically with $N$ (cubic work growth),
    from tens of thousands at $N=128$ to billions (L2) / tens of millions
    (LLC) at $N=3072$. Across distance and hint, L1D and \texttt{sw\_prefetch\_access.any}
    counts are nearly identical (prefetching does not change the number of
    loads issued, only when the data for them arrives); L2 and LLC misses
    vary more and are generally lower at distance$=64$/hint$=$T1 than at
    other settings, consistent with those prefetches landing closer to the
    actual point of use instead of being re-evicted first.

    \item \textbf{Effect of hardware prefetchers:} Disabling them hurts
    \texttt{naive} far more than \texttt{matmul\_prefetch} at $N=2048$
    ($19$--$44\%$ slower vs.\ $\sim\!10\%$ slower), so the naive/prefetch
    speedup ratio goes \emph{up} when HW prefetch is off ($1.35\times \to
    1.46\times$). This is the clearest confirmation that software prefetch
    and the hardware prefetcher are substitutes for the same job (hiding
    memory latency): \texttt{matmul\_prefetch} has its own mechanism to fall
    back on when the hardware one is removed, while \texttt{naive} has
    nothing, so it absorbs almost the full cost. LLC misses jump
    $\sim\!250\times$ for both kernels when HW prefetch is disabled
    ($2.9$M$\to$$731$M naive, $5$M$\to$$729$M prefetch), confirming the
    hardware prefetcher was previously intercepting the large majority of L2
    misses before they reached the LLC/DRAM.
\end{enumerate}

\section{Task 2B: SIMD (AVX2)}

\subsection{Implementation}

The kernel keeps the natural $i\to j\to p$ loop order but register-blocks the
$j$ (column-of-$B$) loop by \texttt{step}${}=4$: each iteration of the inner
loop computes four output elements $C[i][j\,..\,j{+}3]$ at once. Four
\texttt{\_\_m256} accumulators (\texttt{acc0..acc3}) are held in registers and
zeroed at the top of the block. The $p$ loop then advances 8 floats at a time:
one \texttt{\_mm256\_loadu\_ps} pulls a lane of $A[i][p\,..\,p{+}7]$, which is
reused across all four columns, and four \texttt{\_mm256\_fmadd\_ps} multiply it
by the four $B$-row lanes and add into the accumulators. This $1\times4$
register block amortizes the single $A$ load and the loop overhead over four
FMAs, and the four independent accumulator chains hide the FMA latency.

After the vector loop, each accumulator is reduced to a scalar by
\texttt{hsum256} (an \texttt{extractf128} + \texttt{add} to fold the two
128-bit halves, then two \texttt{hadd\_ps}). Remainders are handled by scalar
cleanup: a \texttt{p < K} tail loop finishes the last $K \bmod 8$ elements into
the same four sums, and a \texttt{j < N} tail loop handles the last
$N \bmod 4$ columns one at a time via the scalar-tail helper \texttt{dot\_simd}.
All loads are unaligned (\texttt{loadu}) so arbitrary \texttt{lda/ldb/ldc} are safe.

\subsection{Baseline Profiling / Instruction Count}

User-space instructions were attributed to each kernel with
\texttt{perf record -e instructions:u} followed by \texttt{perf report}: the
per-call figure is the run's total event count times the kernel's overhead
percentage, divided by the number of times the harness invokes that kernel in
a single-stage run (naive: 1 reference + 1 warmup + 1 + 3 timed ${}={}$ 6;
SIMD: 1 warmup + 1 + 3 ${}={}$ 5). Table~\ref{tab:t2-simd-icount} gives the
$1024^3$ point; Fig.~\ref{fig:t2-simd-icount-vs-size} shows the full size
sweep.

The reduction is essentially size-independent, $\approx 11.8\times$ for the
128-bit implementation and $\approx 23\times$ for the 256-bit implementation across the whole
sweep, because it comes from the code's structure, not from cache
behaviour. It exceeds the raw lane counts (4 and 8) because each vector
iteration also folds a multiply and an add into one FMA and amortizes one
$A$ load, the loop-counter arithmetic and the branch over four columns of
output. Doubling the vector width from 128 to 256 bits roughly halves the
count, as expected.

\begin{table}[H]
    \centering
    \begin{tabular}{lrr}
        \toprule
        \textbf{Stage} & \textbf{Instr./call ($1024^3$)} & \textbf{Reduction vs naive} \\
        \midrule
        \texttt{matmul\_naive}          & $6.45\times10^{9}$ & 1.00$\times$ \\
        \texttt{matmul\_simd} (128-bit)  & $5.45\times10^{8}$ & 11.8$\times$ \\
        \texttt{matmul\_simd} (256-bit)  & $2.79\times10^{8}$ & 23.1$\times$ \\
        \bottomrule
    \end{tabular}
    \caption{User-space instructions per call at $M{=}N{=}K{=}1024$,
    from \texttt{perf record}/\texttt{perf report}.}
    \label{tab:t2-simd-icount}
\end{table}

\begin{figure}[H]
    \centering
    \includegraphics[width=0.95\textwidth]{task2_simd_icount_vs_size.png}
    \caption{Instructions executed per call vs.\ matrix size
    ($M{=}N{=}K$), for \texttt{matmul\_naive}, the 128-bit SIMD, and the
    256-bit implementations. Per-call figures are the per-symbol instruction count
    from \texttt{perf record} (\,\texttt{event\_count} $\times$ overhead\,\%\,)
    divided by the number of times the main function calls each function in a
    single-stage run (naive: 6, SIMD: 5). Log scale.}
    \label{fig:t2-simd-icount-vs-size}
\end{figure}

The 128-bit and 256-bit SIMD bars are absent at $N{=}64$: that workload runs
in a few microseconds and \texttt{perf} collected too few samples so those are ditched here.

\subsection{Performance}

\begin{table}[H]
    \centering
    \begin{tabular}{lrrr}
        \toprule
        \textbf{Stage} & \textbf{Time (ms)} & \textbf{GFLOP/s} & \textbf{Speedup} \\
        \midrule
        \texttt{matmul\_naive} (ref) & 496.506 & 4.33 & 1.00$\times$ \\
        \texttt{matmul\_simd} (256-bit) & 92.051 & 23.33 & 5.39$\times$ \\
        \bottomrule
    \end{tabular}
    \caption{\texttt{./bin/matmul simd 1024 1024 1024 1234}. At smaller, cache-resident sizes the SIMD
    speedup is much higher (up to 19$\times$ at $N{=}256$); see
    Fig.~\ref{fig:t2-simd-width-vs-size}.}
    \label{tab:t2-simd-speedup}
\end{table}

$5.39\times$ is well below the $8\times$ raw lane count of SIMD because at
$N{=}1024$ the $N \times K$ float working set no longer fits in L2/L3: the
kernel is already partly memory-bound here, so the $23\times$ instruction-count
reduction (Table~\ref{tab:t2-simd-icount}) does not translate one-to-one into
performance time speedup.

\subsection{Speedup vs. SIMD Width}

The final \texttt{src/matmul\_simd.cpp} is the 256-bit SIMD implementation. Both implementations were swept over
$M{=}N{=}K \in \{64,\dots,3072\}$ against \texttt{matmul\_naive}; the
speedups are plotted in Fig.~\ref{fig:t2-simd-width-vs-size} and the absolute
runtimes in Fig.~\ref{fig:t2-simd-time-vs-size}. Below $N \approx 512$ the
problem is cache-resident and the wider register roughly doubles the speedup
(19$\times$ vs 8$\times$ at $N{=}256$); beyond the L2/L3 capacity both kernels
become memory-bound and converge to $\approx 4$--$5\times$, where vector width
no longer matters.

\begin{figure}[H]
    \centering
    \includegraphics[width=0.95\textwidth]{task2_simd_speedup_vs_size.png}
    \caption{Speedup of \texttt{matmul\_simd} over \texttt{matmul\_naive} vs.\
    matrix size ($M{=}N{=}K$), for the 128-bit and 256-bit SIMD
    kernels.}
    \label{fig:t2-simd-width-vs-size}
\end{figure}

\begin{figure}[H]
    \centering
    \includegraphics[width=0.95\textwidth]{task2_simd_time_vs_size.png}
    \caption{Kernel execution time (ms, log scale) vs.\ matrix size
    ($M{=}N{=}K$). Same runs as
    Fig.~\ref{fig:t2-simd-width-vs-size}.}
    \label{fig:t2-simd-time-vs-size}
\end{figure}

\begin{table}[H]
    \centering
    \scriptsize
    \resizebox{\textwidth}{!}{%
    \begin{tabular}{ll ccccccccccc}
        \toprule
        & & \multicolumn{11}{c}{\textbf{Matrix size} $N$ ($M{=}N{=}K$)} \\
        \cmidrule(lr){3-13}
        \textbf{SIMD width} & \textbf{Metric} & 64 & 128 & 192 & 256 & 384 & 512 & 768 & 1024 & 1536 & 2048 & 3072 \\
        \midrule
        \multirow{2}{*}{No SIMD (\texttt{matmul\_naive})}
          & Instructions/call & $1.52{\times}10^{6}$ & $1.25{\times}10^{7}$ & $4.31{\times}10^{7}$ & $1.02{\times}10^{8}$ & $3.41{\times}10^{8}$ & $8.08{\times}10^{8}$ & $2.72{\times}10^{9}$ & $6.45{\times}10^{9}$ & $2.18{\times}10^{10}$ & $5.16{\times}10^{10}$ & $1.74{\times}10^{11}$ \\
          & Time (ms) & 0.177 & 1.778 & 8.598 & 8.934 & 21.990 & 59.789 & 201.384 & 496.506 & 1707.046 & 4249.303 & 14687.692 \\
        \addlinespace
        \multirow{3}{*}{128-bit (SSE)}
          & Instructions/call & --- & $1.15{\times}10^{6}$ & $3.20{\times}10^{6}$ & $9.07{\times}10^{6}$ & $2.92{\times}10^{7}$ & $6.93{\times}10^{7}$ & $2.31{\times}10^{8}$ & $5.45{\times}10^{8}$ & $1.83{\times}10^{9}$ & $4.32{\times}10^{9}$ & $1.46{\times}10^{10}$ \\
          & Time (ms) & 0.056 & 0.395 & 0.844 & 0.843 & 3.362 & 9.502 & 40.829 & 93.164 & 312.891 & 854.104 & 3699.584 \\
          & Speedup vs.\ No SIMD & 6.03$\times$ & 5.51$\times$ & 8.57$\times$ & 8.04$\times$ & 6.61$\times$ & 6.38$\times$ & 5.20$\times$ & 5.26$\times$ & 5.33$\times$ & 4.84$\times$ & 3.98$\times$ \\
        \addlinespace
        \multirow{3}{*}{256-bit (AVX2)}
          & Instructions/call & --- & $6.73{\times}10^{5}$ & $1.77{\times}10^{6}$ & $5.17{\times}10^{6}$ & $1.62{\times}10^{7}$ & $3.61{\times}10^{7}$ & $1.19{\times}10^{8}$ & $2.79{\times}10^{8}$ & $9.30{\times}10^{8}$ & $2.19{\times}10^{9}$ & $7.34{\times}10^{9}$ \\
          & Time (ms) & 0.022 & 0.162 & 0.510 & 0.468 & 1.858 & 6.075 & 38.203 & 92.051 & 354.612 & 843.911 & 3590.953 \\
          & Speedup vs.\ No SIMD & 7.95$\times$ & 10.94$\times$ & 16.88$\times$ & 19.08$\times$ & 11.84$\times$ & 9.84$\times$ & 5.27$\times$ & 5.39$\times$ & 4.81$\times$ & 5.04$\times$ & 4.09$\times$ \\
        \bottomrule
    \end{tabular}}
    \caption{Instructions/call, kernel time, and speedup vs.\ matrix size, for each SIMD
    width. Instructions/call are retired user-space instructions from
    \texttt{perf record}/\texttt{perf report} (per-call = event count $\times$ overhead\%
    divided by function call count); times are the code's median
    function-only \texttt{time(ms)}, not whole process \texttt{perf stat} time. Same source data as
    Table~\ref{tab:t2-simd-icount}, Fig.~\ref{fig:t2-simd-icount-vs-size},
    Fig.~\ref{fig:t2-simd-width-vs-size}, and Fig.~\ref{fig:t2-simd-time-vs-size}.}
    \label{tab:t2-simd-summary}
\end{table}

\subsection{Questions}
\begin{enumerate}[leftmargin=1.6em, itemsep=0.3em]
    \item \textbf{Instruction count change (naive $\to$ SIMD):}
    Instructions per function call drop by about $11.8\times$ for the
    128-bit implementation and $23\times$ for the 256-bit implementation, roughly constant
    across matrix size (Table~\ref{tab:t2-simd-icount},
    Fig.~\ref{fig:t2-simd-icount-vs-size}). This exceeds the raw lane counts
    (4 and 8) for two reasons: each \texttt{fmadd} replaces a separate
    multiply and add, and the $1\times4$ register block issues one $A$ load
    plus one set of loop-counter/branch instructions per four FMAs instead of
    per one. The count halves cleanly from 128 to 256 bits because the work
    per instruction doubles while the loop structure is identical.

    \item \textbf{Trends across matrix size $\times$ SIMD width; which width
    gives max speedup:}
    For cache-resident sizes ($N \lesssim 512$) speedup rises with $N$ and
    scales with vector width, peaking at $N{=}256$: $19.1\times$ for 256-bit
    and $8.0\times$ for 128-bit. Here the code is compute-bound, so halving
    the instruction count with a wider register directly halves the runtime.
    As $N$ grows past the L2/L3 capacity the $N\times K$ floats
    no longer fits and the function streams it from memory on every
    row of $A$; both widths become memory-bound and converge to
    $\approx 4$--$5\times$, where SIMD width stops mattering. So 256-bit
    gives the maximum speedup at every size, but its advantage over
    128-bit shrinks from $\sim 2\times$ to negligible as the workload leaves
    cache.
\end{enumerate}


\section{Task 2C: Software Prefetching + SIMD (Combined / Optimized)}

\subsection{Implementation}

\texttt{matmul\_optimized} tiles $N$ into blocks of $T=96$ columns of $B$. Within a
tile it uses $4\times3$ SIMD "array" (4 rows of $A$
$\times$ 3 columns of $B$, 12 \texttt{\_\_m256} accumulators) with two
manually-unrolled $K$-steps of 8 floats each (16 floats/iteration) before a
scalar tail. Software prefetch (\texttt{\_MM\_HINT\_T0}, distance $64$
floats, same as Task 2A) is issued on the 3 active $B$ rows only, since $B$ is
re-streamed once per internal loop iteration while the 4 $A$ rows are reused across
all $j$ in a tile and mostly stay resident. Arbitrary $M/N/K/lda/ldb/ldc$ are
handled with scalar fallback loops on every edge: a $1$-row tail when $M \bmod
4 \ne 0$, a $1$-column tail when a $T$-tile's width isn't a multiple of 3, and
a scalar (with SIMD) dot-product tail on $K \bmod 8$. Block size and unroll
depth were tuned by sweeping $T$ on $N{=}1024,2048$.
$T=96$ (a multiple of the $4\times3$ tile) and the $4\times3$ shape were
chosen because they exactly saturate the 16 available SIMD registers with
no spilling: the inner loop keeps 12 accumulators + 3 live $B$ vectors + 1
reused $A$ vector $= 16$ registers live at once.


\subsection{Final Performance Summary}

$N\in\{128,256,512,1024,1536,2048,3072\}$, seed $1234$. Speedup is each stage
vs.\ \texttt{matmul\_naive} on the same workload.

\begin{figure}[H]
    \centering
    \includegraphics[width=0.75\textwidth]{task2_all_opts_vs_size.png}
    \caption{Speedup of \texttt{matmul\_simd}, \texttt{matmul\_prefetch}, and
    \texttt{matmul\_optimized} over \texttt{matmul\_naive} vs.\ matrix size
    (log-log).}
    \label{fig:t2-all-opts-vs-size}
\end{figure}

\begin{table}[H]
    \centering
    \begin{tabular}{lrr}
        \toprule
        \textbf{Stage} & \textbf{Speedup @ 1024$^3$} & \textbf{Speedup @ 2048$^3$} \\
        \midrule
        \texttt{matmul\_simd}      & $5.90\times$  & $5.79\times$  \\
        \texttt{matmul\_prefetch}  & $1.37\times$  & $1.33\times$  \\
        \texttt{matmul\_optimized} & $23.81\times$ & $22.74\times$ \\
        \bottomrule
    \end{tabular}
    \caption{Speedup of each stage (median of 2 repeats).}
    \label{tab:t2-all-opts-table}
\end{table}

\begin{figure}[H]
    \centering
    \includegraphics[width=0.95\textwidth]{task2c_cache_metrics.png}
    \caption{L1D / L2 / LLC miss counts and \texttt{sw\_prefetch\_access.any}
    vs.\ matrix size, one line per stage (raw \texttt{perf stat} counts, log-log).}
    \label{fig:t2-all-opts-cache}
\end{figure}

\subsection{Final Score}

Full graded run (\texttt{./bin/matmul}, default $N{=}1024$, seed $1234$):

\begin{verbatim}
stage           correct    time(ms)     GFLOP/s    speedup
--------------------------------------------------------------
naive (ref)     yes         518.481        4.14      1.00x
simd            yes          83.772       25.63      6.19x
prefetch        yes         376.951        5.70      1.38x
optimized       yes          22.284       96.37     23.27x
--------------------------------------------------------------
correctness : 30 / 30   (3 of 3 stages correct)
optimized   : 23.27x  ->  55 / 70  speedup points
TOTAL       : 85 / 100
\end{verbatim}

\subsection{Discussion}

SIMD comes before tiling because the naive/scalar loop is compute-bound (a
single dependent FMA chain per $(i,j)$ over $K$): with nothing to overlap,
cache blocking cannot show a benefit until the inner loop is fast enough to
be memory-bound instead. This is visible in \texttt{matmul\_prefetch}
(scalar, tiled, prefetched): it fixes the memory side but stays compute-bound
at scalar throughput, so its speedup is a flat $\sim 1.34\times$ regardless
of size. Tiling and prefetch have nothing left to hide once the loop
itself is the bottleneck.

\texttt{matmul\_simd} shows the opposite failure: fast compute, but no tiling. Its speedup collapses from $\sim 13$--$15\times$ at $N
\le 256$ (fits in L2/L3) to $\sim 4$--$6\times$ at $N \ge 1024$, matching the
cache-metrics plot: L1D and LLC misses for \texttt{matmul\_simd} grow far faster
with $N$ than for \texttt{matmul\_optimized} or \texttt{matmul\_prefetch} (e.g.\
at $N{=}2048$, LLC misses are $101$M for \texttt{simd} vs.\ $4.2$M for
\texttt{optimized} and $1.8$M for \texttt{prefetch}). With no tiling, $B$ is
re-streamed from memory on every row of $A$ once it no longer fits in cache.

\texttt{matmul\_optimized} combines both fixes and inherits the best of each:
SIMD throughput plus \texttt{matmul\_prefetch}-level cache
behavior from the $N$-tiling, so its speedup stays flat at $\sim 22$--$24\times$
for $N \ge 512$ instead of decaying like \texttt{matmul\_simd}. Prefetch adds a
further, smaller margin on top by hiding the residual DRAM latency on the 3
active $B$ rows per internal loop iteration. \texttt{sw\_prefetch\_access.any} is
essentially zero for \texttt{matmul\_simd} (no explicit prefetch) and large for
both \texttt{matmul\_prefetch} and \texttt{matmul\_optimized}, confirming it is
this issued-prefetch traffic, not incidental hardware prefetching, driving the
gap.

\section{llama.cpp Results}


\begin{table}[H]
    \centering
    \begin{tabular}{lrrr}
        \toprule
        \textbf{Metric} & \textbf{llama.cpp native} & \textbf{student kernel} & \textbf{Speedup} \\
        \midrule
        Prompt-eval (tok/s) & 12{,}562.81 & 16{,}129.03 & $1.28\times$ \\
        Generation (tok/s)  & 7{,}437.14  & 9{,}826.86  & $1.32\times$ \\
        \bottomrule
    \end{tabular}
    \caption{\texttt{make llama-demo} comparison.}
    \label{tab:t2-llama}
\end{table}

\end{document}
